\documentclass[a4paper,11pt]{article}

\usepackage[utf8]{inputenc}
\usepackage[T1]{fontenc}
\usepackage{babel}
\babelprovide[import, main]{spanish}
\usepackage[top=1.6cm,bottom=1.6cm,left=1.7cm,right=1.7cm]{geometry}
\usepackage{booktabs}
\usepackage{array}
\usepackage{tabularx}
\usepackage{enumitem}
\usepackage{setspace}
\usepackage{xcolor}
\usepackage{graphicx}
\usepackage[colorlinks=true,linkcolor=black,urlcolor=blue]{hyperref}

\setlength{\parindent}{0pt}
\setlength{\parskip}{0.35em}

\newcommand{\sect}[1]{\vspace{0.4em}\section*{\normalfont\large\bfseries\MakeUppercase{#1}}\vspace{-0.2em}}

\title{Séptima Ola --- Requisitos de Contratación y Logística}
\author{Séptima Ola}
\date{\today}

\begin{document}

\begin{center}
    {\LARGE\bfseries SÉPTIMA OLA}\\[0.25em]
    {\Large Requisitos de Contratación y Logística (Booking)}\\[0.15em]
    {\small Tarifas, Formatos de Presentación, Cobertura y Traslados}\\[0.3em]
    {\small \today}
\end{center}

\vspace{0.2em}
\hrule
\vspace{0.5em}

\sect{Tarifas y Formato del Show}

\begin{itemize}[leftmargin=1.4em, itemsep=0.2em]
    \item \textbf{Tarifa de presentación:} \$3,000 MXN por hora de show (honorarios base).
    \item \textbf{Duración del show regular:} La presentación estándar de Séptima Ola dura aproximadamente \textbf{40 minutos}, con un repertorio compuesto principalmente por temas originales que fusionan reggae, ska, rocksteady y jazz.
    \item \textbf{Extensión y adaptación del show:} En caso de que el evento, festival o recinto requiera una duración mayor a una hora (ej. 60, 90, 120 o más minutos / varios sets), la banda puede ampliar la presentación para cumplir a cabalidad con el tiempo requerido.
    \item \textbf{Versatilidad de repertorio:} Gracias a la amplia experiencia y versatilidad técnica de los integrantes, el set extendido puede incorporar versiones clásicas de reggae y ska, así como covers populares de otros géneros (rock, pop, música latina, jazz) adaptados al estilo característico de la banda, además de arreglos extendidos, solos instrumentales y dinámicas con el público.
\end{itemize}

\sect{Cobertura Geográfica y Traslados}

\begin{itemize}[leftmargin=1.4em, itemsep=0.2em]
    \item \textbf{Base de operaciones:} La Raza, Ciudad de México (CDMX).
    \item \textbf{Zona local (CDMX y Estado de México):} Aplica la tarifa base de presentación. La logística y traslado local del elenco y equipo dentro de la Ciudad de México y los municipios conurbados del Estado de México corren por cuenta de la banda.
    \item \textbf{Presentaciones foráneas (Fuera de CDMX y Edo. de México):} Para fechas fuera de la zona metropolitana, el organizador o promotor debe cubrir los gastos adicionales de transporte, casetas de peaje, viáticos de alimentación y, de requerirse por itinerario, hospedaje.
\end{itemize}

\sect{Costos de Referencia para Transporte Foráneo}

El equipo de viaje de Séptima Ola para presentaciones fuera de la ciudad se compone de \textbf{8 personas} (6 músicos en escenario + 2 integrantes de equipo técnico: Ingeniera de sonido y Fotógrafa), además de backline e instrumentos musicales.

\begin{center}
    \renewcommand{\arraystretch}{1.25}
    \small
    \begin{tabularx}{\textwidth}{@{} >{\raggedright\arraybackslash}p{3.2cm} >{\raggedright\arraybackslash}p{2.5cm} >{\raggedright\arraybackslash}p{2.8cm} >{\raggedright\arraybackslash}X @{}}
        \toprule
        \textbf{Modalidad} & \textbf{Costo Aprox. por Persona} & \textbf{Costo Total Estimado (8 Personas + Equipo)} & \textbf{Detalles e Inclusiones} \\
        \midrule
        \textbf{Van Privada con Chofer} \newline\emph{(Recomendado)} & \$700 -- \$1,500 MXN & \$5,500 -- \$12,000 MXN / viaje & Camioneta de pasajeros (Toyota Hiace o Mercedes Sprinter, 13--15+ plazas) con capacidad para 8 pasajeros y carga de instrumentos; incluye renta con operador (\$3,500--\$6,000/día), combustible, casetas de peaje y viáticos del chofer. \\
        \midrule
        \textbf{Autobús Foráneo de Línea} & \$700 -- \$1,800 MXN \newline(viaje redondo) & \$5,600 -- \$14,400 MXN \newline(viaje redondo) & Boletos de líneas comerciales (ADO, Primera Plus, ETN) para distancias intermedias (150--450~km). Requiere fletes o taxis locales adicionales para el traslado de batería y equipo pesado. \\
        \midrule
        \textbf{Vuelos Nacionales} & \$1,800 -- \$4,000 MXN \newline(viaje redondo) & \$18,000 -- \$38,000+ MXN \newline(total estimado) & Vuelos en aerolíneas nacionales (Volaris, VivaAerobus, Aeroméxico) para destinos de larga distancia (ej. Monterrey, Guadalajara, Tijuana, Cancún) más cargos de equipaje/instrumentos (\$1,000--\$2,000 MXN por tramo). \\
        \bottomrule
    \end{tabularx}
\end{center}

\sect{Viáticos y Hospedaje (En Caso de Aplicar)}

\begin{itemize}[leftmargin=1.4em, itemsep=0.2em]
    \item \textbf{Alimentación / Viáticos:} \$350 a \$500 MXN por persona al día (equivalente a \$2,800 -- \$4,000 MXN diarios para el equipo de 8 personas) cuando la producción o el evento no proporcione servicio directo de catering o alimentos.
    \item \textbf{Hospedaje:} 3 a 4 habitaciones dobles en hotel seguro, cómodo y con estacionamiento, preferentemente ubicado en las inmediaciones del recinto, cuando la distancia o el horario del evento obliguen a pernoctar.
\end{itemize}

\sect{Itinerario y Horarios de Producción}

\begin{itemize}[leftmargin=1.4em, itemsep=0.2em]
    \item \textbf{Llegada y montaje (Load-in):} 90 a 120 minutos antes de la prueba de sonido.
    \item \textbf{Prueba de sonido (Soundcheck):} 45 a 60 minutos previos a la apertura de puertas al público.
    \item \textbf{Desmontaje (Load-out):} 30 a 45 minutos posteriores al término de la presentación.
\end{itemize}

\sect{Contacto para Booking y Contratación}

\begin{itemize}[leftmargin=1.4em, itemsep=0.2em]
    \item \textbf{Correo electrónico:} \href{mailto:septimaolaoficial@gmail.com}{septimaolaoficial@gmail.com}
    \item \textbf{Sitio web oficial:} \href{https://septimaola.com}{https://septimaola.com}
    \item \textbf{Redes sociales:} @septimaolaoficial (Instagram, Facebook, TikTok, YouTube)
    \item \textbf{Ubicación:} La Raza, Ciudad de México, México.
\end{itemize}

\end{document}
